\section{Single-metric relativistic theory}\label{sec:theory}
\subsection{Aether scalar--tensor sector}
Ordinary matter couples minimally and universally to one physical metric $g_{\mu\nu}$. Let $A^\mu$ be a unit timelike vector, $A^\mu A_\mu=-1$, and $\phi$ a shift-symmetric scalar. Define
\begin{align}
F_{\mu\nu}&=2\nabla_{[\mu}A_{\nu]},&
J^\mu&=A^\nu\nabla_\nu A^\mu,\\
q_{\mu\nu}&=g_{\mu\nu}+A_\mu A_\nu,&
\Q&=A^\mu\nabla_\mu\phi,\\
\Y&=q^{\mu\nu}\nabla_\mu\phi\nabla_\nu\phi.
\end{align}
The covariant scalar--vector sector is the AeST action \cite{SkordisZlosnik2021,SkordisZlosnik2022,BatakiSkordisZlosnik2024}, with bare cosmological constant fixed to zero,
\begin{widetext}
\begin{align}
S_{\rm AeST}=&\frac{1}{16\pi\widetilde G}\int d^4x\sqrt{-g}\Big\{R-\frac{K_B}{2}F_{\mu\nu}F^{\mu\nu}
 +(2-K_B)\left(2J^\mu\nabla_\mu\phi-\Y\right)
-\Ffull(\Y,\Q)-\lambda(A^\mu A_\mu+1)\Big\}
+S_m[g_{\mu\nu},\psi_m].
\label{eq:aest-action}
\end{align}
\end{widetext}
No disformal matter metric is introduced.

\subsection{Nonseparable scalar kinetic function}
The scalar free function is
\begin{equation}
\Ffull(\Y,\Q)=\FT(\Q)+\Y\Gamma(\Q)+J(\Y),
\label{eq:fullF}
\end{equation}
where
\begin{equation}
\FT(\Q)=-4K_2Z_0^2\left[\cosh\left(\frac{\Q-\Q_0}{Z_0}\right)-1\right],
\label{eq:FT}
\end{equation}
\begin{equation}
\Gamma(\Q)=\frac{2K_2Z_0}{\Q}\sinh\left(\frac{\Q-\Q_0}{Z_0}\right),
\label{eq:Gamma}
\end{equation}
and the static galactic function is
\begin{align}
J(\Y)=&\frac{a_0^2(2-K_B)}{\beta_0^3}\Bigg[-\frac{2\sqrt{\Y}\beta_0(\beta_0+1)}{a_0}
+\frac{\Y\beta_0^2}{a_0^2}\nonumber\\
&+2(\beta_0+1)^2\ln\left(1+\frac{\sqrt{\Y}\beta_0}{a_0(\beta_0+1)}\right)\Bigg].
\label{eq:JY}
\end{align}
The relevant temporal derivatives are
\begin{align}
\mathcal{F}_{T,\Q}&=-4K_2Z_0\sinh\left(\frac{\Q-\Q_0}{Z_0}\right),\\
\mathcal{F}_{T,\Q\Q}&=-4K_2\cosh\left(\frac{\Q-\Q_0}{Z_0}\right),\\
\mathcal{F}_{T,\Q\Q\Q}&=-\frac{4K_2}{Z_0}\sinh\left(\frac{\Q-\Q_0}{Z_0}\right),
\end{align}
and
\begin{align}
\Gamma_{,\Q}=\frac{2K_2}{\Q^2}\Bigg\{&\Q\cosh\left(\frac{\Q-\Q_0}{Z_0}\right)\nonumber\\
&-Z_0\sinh\left(\frac{\Q-\Q_0}{Z_0}\right)\Bigg\}.
\label{eq:GammaQ}
\end{align}
The essential minimality property is that $\Gamma$ is fixed by the homogeneous current, not independently fitted. Removing the mixed term provides a separable control theory; no additional field or continuous parameter is required to pass from that control to Eq.~\eqref{eq:fullF}.

\subsection{Auxiliary acceleration sector}
Late-time acceleration is supplied by a Type-II minimally modified gravity sector of the VCDM class \cite{DeFeliceMukohyamaPookkillath2021,DeFeliceMaedaMukohyamaPookkillath2022,GanzMartensMukohyamaNamba2023,AkarsuEtAl2024VCDM}. In ADM variables, the part relevant for the auxiliary scalar $\varphi_V$ and its constraint multiplier can be written, up to the standard Einstein--Hilbert ADM contribution and gauge-fixing term, as
\begin{align}
S_V=\mpl^2\!\int dt\,d^3x\,N\sqrt{\gamma}\Big[&-V(\varphi_V)-\frac34\lambda_C^2\nonumber\\
&-\lambda_C(K+\varphi_V)-\mathcal L_{\rm gf}\Big].
\label{eq:vcdm-aux}
\end{align}
Varying $\lambda_C$ gives
\begin{equation}
\lambda_C=-\frac23(K+\varphi_V),
\end{equation}
so the reduced deformation is
\begin{equation}
\mathcal L_{V,\rm red}=\mpl^2\left[\frac{(K+\varphi_V)^2}{3}-V(\varphi_V)-\mathcal L_{\rm gf}\right].
\label{eq:vcdm-red}
\end{equation}
For generic scalar perturbations with $k\neq0$, the auxiliary equations enforce $\delta\varphi_V=0$ together with the corresponding multiplier constraints; after Schur elimination this sector adds no propagating scalar degree of freedom. This is the key reason it can alter the background acceleration history without introducing a dark-energy fluid or an extra scalar wave mode \cite{DeFeliceMaedaMukohyamaPookkillath2022}.

In the numerical background parameterization, the reconstructed acceleration history is represented by a fourth-order basis with coefficients $s_1,\ldots,s_4$ and present closure density $\Omega_{V}$. Flat closure is recomputed at every likelihood point,
\begin{equation}
\Omega_V=1-\frac{\omega_b+\omega_{\rm aest}+\omega_r}{h^2},
\label{eq:vcdm-closure}
\end{equation}
for the zero-particle-CDM realization, while one dependent polynomial coefficient is fixed by the basis normalization. The final fit varies only one independent shape coefficient, $s_4$, from this acceleration sector.

\subsection{Regime separation}
The construction assigns distinct roles to distinct operators. The temporal scalar sector determines the homogeneous effective cold component; $J(\Y)$ determines the quasistatic MOND response; the mixed-current operator controls the radiation-era scalar gradient structure; and the auxiliary Type-II constraint sector supplies the reconstructed late-time acceleration. This separation is used throughout the validation to avoid attributing a successful observable to an unnecessarily complicated sector.
